\chapter{Exposé}

\section{Einleitung}

Der \emph{Deep Space 8K}%
\footnote{\url{https://ars.electronica.art/center/de/exhibitions/deepspace/}}
des Ars Electronica Centers in Linz bietet mit seiner $16 \times 9$~m
großen Projektionsfläche inklusive Positionstracking eine einzigartige
Möglichkeit, Computerspiele zu realisieren. Diese Spiele verwenden keine
klassischen Kontrollmechanismen wie Tastatur, Maus oder Gamepad sondern die
Spieler*innen selbst "steuern" die Inhalte mit ihren Bewegungen. Darüber
hinaus finden diese Spiele in einem halb-öffentlichen bis öffentlichen Raum
statt, wodurch sich die Bestimmung der Zielgruppe sowie die Anzahl der
spielenden Personen schwierig gestaltet. Diese Bachelorarbeit (Masterarbeit)
beleuchtet diese Problematik und stellt konkrete Lösungsvorschläge anhand
eines Beispiels dar.


\section{Theoretischer Hintergrund und Stand der Forschung}
\label{sec:hintergrund}

Large Public Display Games (LPD Games) sind Spiele, die auf großen,
öffentlich einsehbaren Projektionsflächen dargestellt werden. Derlei
Installationen finden sich etwa in Museen (wie dem Ars Electronica Center)
oder auch auf öffentlichen Plätzen. Personen können diese Spiele in der Regel
jederzeit sehen und auch aktiv an ihnen teilnehmen. Durch diese
Öffentlichkeit ergeben sich nach \cite{Finke2008} drei Arten von
Personengruppen, die am Spiel beteiligt sind: \emph{Actors} nehmen aktiv am
Spielgeschehen teil, \emph{Spectators} verfolgen dieses aktiv und
\emph{Bystanders} befinden sich lediglich in der Umgebung der öffentlichen
Installation. Das Ziel ist es, dass Bystanders zur Spectators und Spectators
zu Actors werden, also das Spiel aktiv spielen. Dieser Prozess soll dabei
möglichst fließend vonstattengehen und eine größtmögliche Anzahl an Personen
umfassen. Ein derartiger Ansatz wurde in \cite{Hochleitner2013} als
\emph{Smooth Transition Gameplay} bezeichnet. Anhand einer konkreten
Anwendung wird dabei demonstriert, wie dieser Übergang erreicht werden kann,
es wird jedoch nicht systematisch beschrieben, welche Faktoren dafür nötig sind.

Einen möglichen Ansatzpunkt bieten dabei die verwendeten Spielmechaniken. Der
in \cite{Schell2019} aufgestellten Kategorisierung folgend bieten sich hierbei 
vor allem Mechaniken aus den Bereichen Raum (\emph{Space}), Handlungen
(\emph{Actions}) und Regeln (\emph{Rules}) an. Dort angesiedelte Mechaniken
können in einem entsprechenden Gamedesign so eingesetzt werden, dass in einem
LPD Game die oben genannten Anforderungen -- möglichst einfacher Einstieg und
gute Skalierbarkeit in Bezug auf die Anzahl der Spieler*innen -- erreicht
werden (siehe Abbildung \ref{fig:logo}). %nur ein Test

\begin{figure}
	\centering
	\includegraphics[width=0.3\textwidth]{logo}
	\caption{Diese Abbildung dient nur als Test.}
	\label{fig:logo}
\end{figure}

\section{Forschungsfrage}

Aus diesen Ansätzen ergibt sich die folgende Forschungsfrage für diese
Bachelorarbeit (Masterarbeit):
%
\begin{quote}
	Welche Spielmechaniken müssen auf welche Art und Weise in einem
	Game\-design für ein Large Public Display Game eingesetzt werden, um
	dieses für eine variable Anzahl von Spieler*innen zu gestalten und diesen
	einen möglichst leichten Einstieg zu ermöglichen?
\end{quote}


\section{Methodik}

Um diese Frage zu beantworten, soll die Bachelorarbeit (Masterarbeit) als eine
Kombination von Literaturarbeit und praktischer \bzw prototypischer Umsetzung
realisiert werden.

Zunächst soll aus bestehender Literatur (erweiternd zu Abschnitt
\ref{sec:hintergrund}) erörtert werden, wie mit dem Thema des Smooth
Transition Gameplay aus Sicht des Gamedesigns umgegangen wurde. Gemeinsame
Faktoren wie Mechaniken sollen daraus extrahiert werden und als Grundlage für
ein eigenes, theoretisches Framework dienen. Dieses Framework soll
schlussendlich eine Liste von Kernmechaniken und Richtlinien für deren
Anwendung enthalten, sodass LPD Games einen leichten Einstieg sowie eine
variable Anzahl von Spieler*innen ermöglichen.

Überprüft soll die Anwendbarkeit dieses Frameworks durch ein eigenes, im
Rahmen des Semesterprojekts 5 entwickeltes, LPD Game werden. Durch einfache,
qualitative Fragestellungen an die Spieler*innen und Beobachtungen der
Besucher*innen während mehrerer Testläufe soll herausgefunden werden, ob der
Gedanke des Smooth Transition Gameplays mit den verwendeten Mechaniken
erreicht werden konnte.


\section{Erwartete Ergebnisse}

Als konkretes Ergebnis wird ein Framework aus Spielmechaniken erstellt,
welches als Grundlage für die Erstellung von LPD Games dienen soll. Es wird
erwartet, dass sich solche konkreten Mechaniken finden und beschreiben lassen.

Bei den Tests der praktischen Umsetzung des Frameworks wird ebenfalls eine
positive Evaluierung erwartet, da es bereits erfolgreiche Konzepte \bzw LPD
Games gibt, auf deren Erfahrungen aufgebaut werden kann.
